\documentclass[12pt]{article}
\usepackage[T1]{fontenc}
\usepackage[utf8]{inputenc}
\usepackage[brazil]{babel}
\usepackage[margin=1in]{geometry}
\setlength{\parindent}{0pt}
\begin{document}
\section*{Tema 57}
Processo: PEDILEF 0508032-49.2007.4.05.8201/ PB\\
Situacao: Julgado\\
Ramo: DIREITO PREVIDENCIÁRIO\\
Questao: Saber se há incidência da prescrição do fundo de direito para os casos em que transcorrido o prazo de 5 anos no período compreendido entre o cancelamento administrativo do auxílio-doença e o ajuizamento da ação.\\
Tese: Não ocorre prescrição do fundo de direito quando, entre o cancelamento administrativo do auxílio-doença e o ajuizamento da ação, decorrem mais de 5 anos.\\
Fonte: https://www.cjf.jus.br/phpdoc/virtus/
\end{document}
